\begin{abstract}\label{sec:abstract}
	Las picaduras de mosquitos transmiten enfermedades como dengue, zika y malaria, por lo que \textbf{detectar su actividad de vuelo} no es un fin en sí mismo sino un habilitador de \textbf{acciones de control vectorial} con impacto sanitario. En este trabajo se explicita primero el \emph{para qué}: (i) \textbf{vigilancia continua} y de bajo costo para estimar abundancia, estacionalidad y generar \emph{alertas tempranas}; (ii) \textbf{identificación de especies} a partir de su firma de aleteo para priorizar vectores de mayor riesgo; (iii) \textbf{mapeo espacial} de focos y rutas de vuelo para orientar operativos; y (iv) \textbf{intervención focalizada} (p.,ej., activación de un dron rociador o trampa inteligente) que reduce insecticida y exposición humana. A partir de estos usos, se conectan motivaciones y medios técnicos: se revisan enfoques clásicos (phase-locked loops y filtros casados) útiles para la \emph{detección} bajo supuestos controlados; se incorporan métodos de \emph{dirección de arribo} (beamforming, MUSIC, ESPRIT) que añaden \emph{localización}; y se discuten técnicas modernas de \emph{aprendizaje profundo} (convolucional y bayesiano) que aumentan robustez ante ruido y variabilidad ambiental, incluida su integración con DoA para seguimiento de trayectorias. Finalmente, se consideran \textbf{aceleradores en FPGA} como vía práctica hacia \emph{tiempo real} y despliegue a bordo de drones o nodos IoT con restricciones de energía. El aporte es una \textbf{síntesis orientada a decisiones}: criterios para elegir tecnología según objetivo operativo (vigilancia vs.\ intercepción), condiciones acústicas y presupuesto computacional, mostrando la transición desde DSP clásico a enfoques híbridos para \textbf{detectar, ubicar y actuar} frente a mosquitos en escenarios reales de salud pública.
\end{abstract}

\begin{IEEEkeywords}
	Salud pública, vigilancia vectorial, detección acústica de mosquitos, dirección de arribo, aprendizaje profundo, drones, FPGA.
\end{IEEEkeywords}